\subsection{Arbitrageurs, order anticipators and bluffers}

\begin{qfdefinition}[Arbitrage: hedge portfolio, basis, fair value, arbitrage spread and bounds; pure and speculative arbitrage]\label{def:arbitrage}
\qffield{Statement} Arbitrageurs trade instruments with correlated prices. An \emph{arbitrage opportunity} arises when those prices diverge from their normal relation: the arbitrageur buys the relatively cheap instruments, sells the relatively expensive ones, and profits if the prices converge \citep[chap.~``Arbitrageurs'', sec.~17.1, p.~348]{alpha-harris}. The positions form a \emph{hedge portfolio}, whose legs are scaled by hedge ratios against a reference instrument, the \emph{arbitrage numerator}, chosen to minimise the portfolio's risk. Its \emph{carrying costs} include financing, dividends paid on short positions, storage and spoilage, offset by dividends, lending fees and interest received \citep[same sec., pp.~348--349]{alpha-harris}. With $B_t$ the \emph{basis}, the difference in prices between the instruments in the hedge portfolio, and $B_t^\star$ its \emph{fair value}, the basis if all were correctly priced relative to each other, which depends on carrying costs and must be estimated, the \emph{arbitrage spread} is
\[ A_t=B_t-B_t^\star \]
\citep[same sec., p.~349]{alpha-harris}. Arbitrageurs trade only when the basis lies outside the \emph{arbitrage bounds}, the values on either side of fair value at which they are just willing to trade \citep[same sec., p.~349]{alpha-harris}. Interpretation: if one unit of basis converges fully to fair value, the trade earns $|A_t|$ gross; with all-in round-trip cost $K_t$ and a charge $R_t$ for residual basis risk it is attractive only when $|A_t|>K_t+R_t$, so the bounds are $B_t^\star\pm(K_t+R_t)$.

\qffield{Conditions} Timeless (concepts); examples to 2000. Harris distinguishes two types by the behaviour of the hedge portfolio's value \citep[same chap., sec.~17.3, pp.~350--363]{alpha-harris}:
\begin{itemize}
\item \emph{pure arbitrage}: the value is strictly mean reverting, because some mechanism forces convergence. The mechanisms are \emph{shipping} (the same instrument in two places, bounded by shipping and financing costs), \emph{delivery} (contract and cash prices converging at maturity, the basis set by carrying costs before it) and \emph{conversion} (the same risk in different forms, bounded by the costs of converting and carrying). Cyclic arbitrages such as currency cross-rates are a kind of conversion arbitrage.

\emph{Cyclic currency arbitrage with quotes}. Write $S_{X/Y}$ for units of $Y$ per unit of $X$. The exact equality $S_{A/B}S_{B/C}=S_{A/C}$ is the no-arbitrage condition only in a frictionless market with one price per pair; with quotes it need not hold for midpoints. For example, with EUR/USD 1.0999/1.1001, USD/JPY 109.99/110.01 and EUR/JPY 120.99/121.03, the synthetic midpoint is 121.00 and the direct one 121.01, yet the two cycles return 0.99957 (via USD) and 0.99974 (via JPY) per euro, so neither is profitable.

With bid and ask, converting $X$ into $Y$ earns the bid of the $X/Y$ quote and converting $Y$ into $X$ pays its ask; starting from one unit of EUR, the two directions of the EUR--USD--JPY cycle return $b_{\text{EUR/USD}}\,b_{\text{USD/JPY}}/a_{\text{EUR/JPY}}$ and $b_{\text{EUR/JPY}}/(a_{\text{USD/JPY}}\,a_{\text{EUR/USD}})$, so there is no cyclic arbitrage if and only if the direct quote overlaps the synthetic band:
\[ b_{\text{EUR/USD}}\,b_{\text{USD/JPY}}\le a_{\text{EUR/JPY}}\quad\text{and}\quad b_{\text{EUR/JPY}}\le a_{\text{EUR/USD}}\,a_{\text{USD/JPY}} \]
Arbitrage parity holds only within this band, widened by spreads, funding and settlement costs; it is not an equality.
\item \emph{speculative arbitrage}, also called risk arbitrage: the value is nonstationary, because instrument-specific factors make it wander, and arbitrageurs bet that short-run mean reversion dominates. Kinds: spreads (maturity, calendar, yield-curve and credit spreads), pairs trading, statistical arbitrage with factor models, and merger arbitrage, the usual sense of ``risk arbitrage''.
\end{itemize}
Arbitrage is understood two ways \citep[same chap., secs.~17.6 and~17.9, pp.~375--378]{alpha-harris}. In the \emph{price characterisation}, arbitrageurs are informed traders in relative values who enforce the law of one price \citep[chap.~``Informed Traders and Market Efficiency'', sec.~10.5.4, pp.~232--235]{alpha-harris}. In the \emph{quantity characterisation}, they discipline slow prices, ship common-factor risk between markets and repackage it, manufacturing derivative positions; the arbitrage spread is the payment for these services.

\qffield{Practice} Arbitrageurs are dealers or value traders in the hedge portfolio: fast dealer-type arbitrageurs trade short-term reversals, and value-type arbitrageurs take large positions when sure of fair value \citep[chap.~``Arbitrageurs'', sec.~17.2, pp.~349--350]{alpha-harris}. Dealers span time, matching buyers and sellers who arrive at different times in one market; arbitrageurs span place, matching those who arrive together in different markets \citep[same chap., sec.~17.7, p.~376]{alpha-harris}. Free entry narrows arbitrage spreads until they just cover the costs of the most efficient arbitrageurs. These are often traders with a natural interest in one leg, such as grain operators near a delivery point \citep[same chap., sec.~17.8, pp.~377--378]{alpha-harris}.

\emph{Economic function}. Arbitrageurs ``buy in cheaper markets and sell in more expensive markets'', connecting sellers in one market to buyers in another \citep[chap.~1, p.~6]{alpha-harris}; as liquidity suppliers they are ``not suppliers of liquidity but porters of liquidity'', demanding it where it is available and supplying it where it is demanded, and they compete with dealers, who connect buyers and sellers arriving at different times in one market \citep[chap.~``Liquidity'', sec.~19.3.6, pp.~404--405]{alpha-harris}. They enforce the law of one price only when costs, funding, hedge error and execution risk are small enough (\hyperref[res:arbitrage-risks]{arbitrage risks}).

\emph{Cross-asset and cross-venue arbitrage}. Dealers connect buyers and sellers across time, while arbitrageurs connect them across instruments, venues, locations and risk forms, often selling a transformed exposure: Harris's dealers span time and arbitrageurs span place, with conversion arbitrageurs as manufacturers: ``the conversion arbitrageur is a manufacturer who produces financial instruments that interest traders'' (sec.~17.3.1.3, p.~354).

Common cross-asset arbitrages map onto Harris's types: cash against futures is a delivery arbitrage; ETF shares against their basket, depositary receipts against local shares and one currency against a synthetic cross rate are conversion arbitrages, the last a cyclic one; option conversions manufacture option exposure from stock and cash (Harris's box ``Manufacturing Stock Option Contracts'', p.~354: arbitrageurs who sell options ``convert equity risk from stock shares to option contracts'' and must rebalance because option values are nonlinear in the stock, pp.~354--355); and merger spreads are speculative risk arbitrage (pp.~360--363).

Harris's arbitrage chapter does not discuss ETFs or depositary receipts; their placement is Interpretation, and ETF creation and redemption is in \hyperref[mech:etf-creation-redemption]{ETF creation and redemption}. \emph{Derivative conversion.} A market maker who buys an option and delta-hedges converts option demand into stock trading, funding and volatility exposure: ``Such hedges are dynamic hedges. They arise whenever the value of one instrument is not an exact linear function of the value of another instrument'', and volatility makes them perform poorly, so option dealers charge for expected hedging costs and manage the risk of their combined positions rather than each arbitrage separately (pp.~354--355). The changing sensitivities include delta, gamma, vega and funding costs; Harris names only delta (n.~1, p.~355).

\qffield{Pitfalls} Full convergence is the assumption behind the $|A_t|$ in the condition above; with partial or slow convergence the gross gain is the expected narrowing net of extra carry (Interpretation; see \hyperref[res:arbitrage-risks]{arbitrage risks}). Pure, speculative, risk and statistical arbitrage are not four parallel types: in Harris risk arbitrage is another name for speculative arbitrage, and statistical arbitrage is one of its kinds. Cross-rate arbitrage is not a fourth pure mechanism beside shipping, delivery and conversion; Harris treats it as a cyclic conversion arbitrage.

Harris's box ``A Three-legged Currency Money Machine'' (p.~357) is internally inconsistent: 110 yen per dollar and 1.10 dollars per euro imply 121 yen per euro, not the printed 100; and its cycle (buy \$1.10 million for 100 million yen, buy 1 million euros for the dollars, buy 99.9 million yen for the euros) loses 0.1 million yen on its own prices, because it buys euros at the implied rate and sells them where they are cheap. A consistent version: with the euro quoted directly at 120.9 yen against the implied 121, buy 1 million euros for 120.9 million yen, sell them for \$1.10 million and the dollars for 121 million yen, a gain of 0.1 million yen, about \$909.

The point of the box, that price pressure from the arbitrage closes the gap and that dealers, the fastest traders, are the natural arbitrageurs, is unaffected. Even pure arbitrage can lose if convergence is slow or financing fails \citep[same chap., secs.~17.3.1 and~17.9, pp.~350 and~378]{alpha-harris}.

\qffield{See also} \hyperref[res:arbitrage-risks]{arbitrage risks}; \hyperref[def:informed-trader-types]{informed-trader types}; \hyperref[def:outside-spread]{outside spread} (value traders' no-trade band, the analogue of arbitrage bounds); \hyperref[def:dealer-inventory-risk]{dealer inventory risk} (dealers span time).

\end{qfdefinition}

\begin{qfresult}[Arbitrage risks: implementation, basis, model and carrying-cost risk, and staying power]\label{res:arbitrage-risks}
\qffield{Statement} Harris names four arbitrage risks \citep[chap.~``Arbitrageurs'', sec.~17.4, pp.~363--373]{alpha-harris}:
\begin{itemize}
\item \emph{implementation risk}: transaction costs larger than expected, from adverse prices on market orders or failed limit orders. It is worst when one leg fills and another does not, when markets move fast, and when the arbitrage spread is small relative to short-term volatility;
\item \emph{basis risk}: the basis moves against the position, because specific factors diverge or common factors are mispriced. Total basis risk scales with the size of the position, and adverse moves force liquidation, at the worst time in a pure arbitrage, when the trader cannot finance the losses;
\item \emph{model risk}: misunderstanding the value relation, so that the trader takes the wrong side, carries unrecognised risks or uses wrong hedge ratios. The commonest error is treating a move of the basis towards a changed fair value as an opportunity;
\item \emph{carrying-cost risk}: higher financing, borrowing, storage or dividend costs than expected; unexpected price rises whose losses on short legs must be funded at once; slow convergence; and forced buy-ins of borrowed securities, which tend to come when the basis has moved against the arbitrageur.
\end{itemize}
Risky arbitrage needs compensation, so arbitrage bounds are wide when arbitrage is risky \citep[same sec., p.~364]{alpha-harris}.

\qffield{Conditions} Timeless (concepts); examples to 2000. Arbitrage opportunities have two causes \citep[same chap., sec.~17.5, pp.~373--375]{alpha-harris}. In the \emph{slow price adjustment} scenario a common factor changes and some prices lag. Arbitrage then moves prices to their new values, and the arbitrageur should trade quickly, demanding liquidity. In the \emph{uninformed liquidity demand} scenario values are unchanged and uninformed traders push prices apart in different markets. The arbitrageur then supplies liquidity, competing with value traders and dealers, and need not trade aggressively. Opportunities do not arise when an instrument-specific value changes: then the fair value of the basis has moved, and a pairs trade against it loses \citep[same chap., pp.~358--360 and~368]{alpha-harris}.

\qffield{Practice} Never leverage to the limit of capital: those who do have no staying power when the basis moves against them \citep[same chap., secs.~17.4.2 and~17.4.4.2, pp.~368 and~372]{alpha-harris}. LTCM, highly leveraged in mean-reverting positions, failed in 1998 although many positions later paid \citep[same chap., box, p.~369]{alpha-harris}. In illiquid markets, trade the illiquid leg first, or interleave small orders in both legs, so as never to be far from hedged \citep[same chap., pp.~366--367]{alpha-harris}. To judge whether a move in the spread is an opportunity, ask which leg moved and why. In Harris's GE against S\&P~500 futures example, a widening caused by the index falling while GE lags is more likely an opportunity than one caused by GE rising, which may be firm-specific news \citep[same chap., box and Table~17-2, pp.~361--362]{alpha-harris}.

\emph{Legging and execution}. Arbitrageurs ``get into trouble when they are stuck in the middle with one leg on and one leg off'' (box, p.~366); they ``try very hard to trade all legs of their hedge portfolio at the same time'', using only market orders when spreads are small relative to the execution uncertainty of limit orders, trading the illiquid leg first when one leg is illiquid so that they can cancel cheaply if the basis closes, and interleaving split orders when both are, often with limit orders (p.~366).

In execution, the arbitrageur estimates the fair basis, costs and hedge ratios, chooses the leg order, submits synchronised orders or trades the illiquid leg first, monitors partial fills, basis and residual exposure, and cancels, completes or rebalances. Trading the illiquid instrument first is Harris's rule; the rest is an illustrative procedure (Interpretation).

\emph{Capital.} Forced liquidation by arbitrageurs who cannot finance losses, meet margin, borrow or carry inventory can widen the basis ``precisely when the forward-looking trade is most attractive'': this is Harris's staying power and carrying-cost risk. That the widest spreads appear when balance sheet is scarce is a generalisation, not a measured result.

\emph{Relative-value features}. Useful features map onto this entry's risks: hedge quality from rolling betas, cointegration or factor residuals and hedge-ratio instability onto model risk; a convergence clock to delivery, expiry, fixing, rebalance or event, and carry from funding, borrow, dividends, storage and roll, onto carrying-cost risk; execution synchrony across legs onto implementation risk; and stress capacity from margin, drawdown to liquidation and recall risk onto staying power. Basis level, carry-adjusted spread, z-score and distance to the bounds measure the opportunity itself. That such features separate cheapness from tradability is an untested interpretation.

\qffield{Pitfalls} Staying power and specific-factor risk are not separate risk categories: in Harris staying power is the consequence of basis and carrying-cost risk under leverage, and specific factors are a source of basis risk and model error. Separating common from specific factors tells whether a spread move is an opportunity or new information, as the passages support. Writing $r_i=\beta_i^\top f+\epsilon_i$ and $r_j=\beta_j^\top f+\epsilon_j$ and attributing spread moves to common factors, a hedge-ratio error $\beta_i-\beta_j$, or residuals $\epsilon_i-\epsilon_j$ holds only for a one-for-one spread.

For a hedge portfolio $r_i-h\,r_j$ with hedge ratio $h$, the spread is $(\beta_i-h\beta_j)^\top f+(\epsilon_i-h\epsilon_j)$, so the factor exposure is $\beta_i-h\beta_j$ and the residual $\epsilon_i-h\epsilon_j$; a hedge ratio chosen to neutralise one factor ($h=\beta_{i1}/\beta_{j1}$) leaves the others (with two factors, $h=1.5$ zeroes the first exposure and leaves $-0.45$ on the second; for independent residuals the spread variance is $(\beta_i-h\beta_j)^\top\Sigma_f(\beta_i-h\beta_j)+\sigma^2_{\epsilon_i}+h^2\sigma^2_{\epsilon_j}$).

\qffield{See also} \hyperref[def:arbitrage]{arbitrage}; \hyperref[def:order-anticipation]{order anticipation} (short squeezes hit arbitrageurs); \hyperref[mech:bubble-crash-dynamics]{bubbles and crashes} (limits to arbitrage).

\end{qfresult}

\begin{qfdefinition}[Order anticipation: front running, quote matching, sentiment-oriented technical trading, squeezes and stop gunning]\label{def:order-anticipation}
\qffield{Statement} \emph{Order anticipators} profit from information about other traders' orders, not about fundamental values, by trading before them \citep[chap.~``Order Anticipators'', pp.~245 and~256--257]{alpha-harris}:
\begin{itemize}
\item \emph{front runners} trade on orders others have decided to arrange. Ahead of aggressive traders they profit from the orders' price impact. Ahead of passive traders they are \emph{quote matchers}, or penny jumpers, who step in front of standing limit orders and keep the limit order as an option to exit: long stock plus the implicit put of a large buy order is a call \citep[sec.~11.1, pp.~245--250]{alpha-harris};
\item \emph{sentiment-oriented technical traders} predict trades uninformed traders have not yet decided to make, from cash-flow timing, hedging needs, asset exchanges, gamblers' interests and speculators' mistakes \citep[sec.~11.2, pp.~251--254]{alpha-harris};
\item \emph{squeezers} corner one side of a market so that those who must close positions have to trade with them \citep[sec.~11.3, pp.~254--256]{alpha-harris}.
\end{itemize}
\emph{Gunning the market}, pushing prices to trigger clustered stop orders and trading with them, is a manipulation of stop orders that Harris treats alongside these strategies \citep[sec.~11.4, p.~256]{alpha-harris}.

\qffield{Conditions} Effects \citep[same chap., secs.~11.1.3--11.1.4, 11.2.1 and~11.3, pp.~250--256]{alpha-harris}. Front running accelerates the impact of the traders it front-runs. It makes prices more informative in the short run when it front-runs informed traders, but in the long run it lowers their profits and investment in information, and so informativeness. It makes prices less informative when it front-runs uninformed traders. It generally makes markets less liquid: front runners' profits are other traders' costs, they offer liquidity back at worse prices, and they drive away the orders they feed on.

Front runners benefit large traders only in narrow conditions: better traders who do not compete with each other and do not squeeze. Sentiment-oriented technical trading makes prices less informative, especially in bubbles, and loses when value traders push prices back, so it pays best in hard-to-value instruments. Squeezes make prices less informative and greatly raise others' costs.

\qffield{Practice} Time precedence with a meaningful minimum price increment makes quote matchers pay to step ahead, to the traders on the other side rather than the passive traders they front-run; a small tick weakens this protection. Passive traders defend themselves by hiding or splitting orders or switching to market orders, which raises their costs and reduces displayed size \citep[same chap., pp.~250 and~257]{alpha-harris}. Traders protect themselves from front runners by guarding order information, from sentiment traders by trading quickly, and from squeezers by keeping several ways to close \citep[same chap., p.~257]{alpha-harris}. NYSE and AMEX specialists, as Harris describes them in 2002, had privileged access to the order flow and used it in the same ways \citep[chap.~``Specialists'', sec.~24.5, pp.~501--508]{alpha-harris}:
\begin{itemize}
\item speculating on short-term forecasts from the order flow;
\item quote matching against a heavy side of their book;
\item \emph{cream skimming}, stepping in front of the book to fill only the more profitable, usually small retail, marketable orders;
\item anticipating stop orders in their books;
\item \emph{stopping stock}, guaranteeing an execution price and then deciding when to fill, a look-back timing option that harms the limit-order traders who would have filled the order.
\end{itemize}
Public order precedence, the public liquidity preservation principle and audit trails limit these privileges \citep[sec.~24.2.2, pp.~498--499]{alpha-harris}.

\qffield{Pitfalls} Not all front running is illegal: it is illegal when the information is obtained improperly, through a confidential brokerage relationship or eavesdropping, but an observant floor trader may legally infer a broker's order \citep[chap.~``Order Anticipators'', sec.~11.1.1, pp.~246--247]{alpha-harris}. Squeezers and gunners defend themselves as speculators, which makes the laws against them hard to enforce \citep[same chap., pp.~256--257]{alpha-harris}. Quote matching and stop gunning are not categories parallel to front running: in Harris quote matching is front running of passive traders, and gunning a manipulation of stop orders. The specialist strategies describe a floor-based market structure of 2002 and are dated.

\qffield{See also} \hyperref[def:bluffing-manipulation]{bluffing and manipulation}; \hyperref[def:informed-trader-types]{informed-trader types} (information-oriented against sentiment-oriented technical trading); \hyperref[res:arbitrage-risks]{arbitrage risks} (short squeezes); \hyperref[mech:bubble-crash-dynamics]{bubbles and crashes} (order anticipators accelerate the rise); \hyperref[res:spread-determinants]{spread determinants} (tick size); \hyperref[def:hft-strategies]{HFT strategies} (order anticipation as one of the SEC's 2010 strategy types, and pinging as a tool).

\end{qfdefinition}

\begin{qfdefinition}[Bluffing and market manipulation: rumourmongers, price manipulators and the asymmetric-impact condition]\label{def:bluffing-manipulation}
\qffield{Statement} \emph{Bluffers} try to fool other traders into trading unwisely and then profit from them. Harris names two techniques \citep[chap.~``Bluffers and Market Manipulation'', p.~259 and Table~12-3, p.~274]{alpha-harris}:
\begin{itemize}
\item \emph{rumourmongers} spread false information, or true information presented so that it will be misread;
\item \emph{price manipulators} arrange trades at prices, volumes and times meant to change others' opinions of value. The trades may be genuine arm's-length trades or \emph{wash trades} with confederates that create artificial activity.
\end{itemize}
\emph{Market manipulation} occurs when bluffers or their victims change prices from what they would otherwise be. Bluffing is profitable when the price impact of the bluffer's purchases differs from that of its sales: if buys move prices more, buy first and sell later, and conversely \citep[same chap., sec.~12.3, pp.~270--273]{alpha-harris}. In Harris's example, liquidity suppliers raise prices 10 cents per 100 contracts for 1{,}000-contract orders but only 5 cents for 500-contract orders. Buying 4{,}000 contracts in 1{,}000s and selling them in 500s then averages 12.50 bought against 12.875 sold, a profit of \$1.5 million on a 1{,}000-unit contract, and leaves the price at 12 rather than 10 \citep[same sec., p.~270]{alpha-harris}. With the same impact per unit for all trades, the same round trip loses \$3 million (p.~271; both figures rechecked, Interpretation).

\qffield{Conditions} Timeless (concepts); legal remarks as of 2002. Bluffers imitate informed speculators: both promote their positions and trade to show confidence. The difference is that speculators trade on well-founded opinions about value, whereas bluffers manufacture the signals; bluffs make prices less informative and are more likely followed by reversals \citep[same chap., sec.~12.2.1, pp.~265--267]{alpha-harris}. Value traders foil bluffs by trading against prices pushed from value, unless bluffers have more capital and force them out. Bluffs therefore succeed most in securities value traders do not follow or cannot trade: illiquid, hard to value or hard to borrow. Long-side bluffs are probably commoner than short-side ones \citep[same chap., sec.~12.2.2, pp.~268--269]{alpha-harris}.

\qffield{Practice} Liquidity suppliers who set prices from net order flow must make buys and sells of every size, speed and timing move prices equally per unit traded; otherwise bluffers can exploit them \citep[same chap., sec.~12.3, pp.~271--273]{alpha-harris}. Momentum traders and liquidity suppliers, who react to the trades they see, are the most vulnerable \citep[same chap., sec.~12.4, p.~274]{alpha-harris}. Transaction costs limit bluffing, which trades intensively \citep[same chap., p.~273]{alpha-harris}.

Symmetric impact disciplines bluffing in this sense. If impact depended on net pressure and risk, of which Harris names only net quantity, a zero-net sequence could not create a higher average sale price than purchase price. That conclusion holds in Harris's example, where each trade's impact is immediate and permanent. In impact models with decay and changing liquidity, equal treatment of buys and sells does not exclude manipulation (see \hyperref[res:no-price-manipulation]{absence of price manipulation}).

\qffield{Pitfalls} A price reversal does not prove manipulation, because good speculators also misjudge values. Conviction usually requires evidence that the bluffer spread information known to be false or arranged wash trades, or testimony on intent; careful bluffers therefore pass as speculators \citep[same chap., sec.~12.2.3, pp.~269--270]{alpha-harris}. Harris's legality column (Table~12-3): false rumours and wash trades are illegal in most jurisdictions and easy to prosecute if caught; misleadingly framed true information and genuine trades meant to mislead are of questionable legality and hard to prove. Wash traders are not a third type beside rumourmongers and price manipulators: in Harris wash trading is a technique of price manipulators.

\qffield{See also} \hyperref[def:order-anticipation]{order anticipation} (gunning stop orders); \hyperref[res:insider-trading-deterrence]{insider-trading detection} (the same proof problem); \hyperref[def:informed-trader-types]{informed-trader types}; \hyperref[mod:kyle-single-auction]{Kyle single-auction model} (Interpretation: impact linear and symmetric in signed flow, as in Harris's no-profit case); \hyperref[res:no-price-manipulation]{absence of price manipulation} (manipulation through one's own impact in a model, not through others' inferences).

\end{qfdefinition}
